\section{Related Work}\label{sec:related}

\subsection{Clinical Diagnosis and Infrared Meibography}

MGD diagnosis combines gland morphology with secretion quality\cite{foulks2003,villani2020latest}. Infrared/near-infrared meibography exploits scattering differences to show dropout, tortuosity, shortening, and cystic dilation\cite{wise2012meibography}. Meiboscore grades dropout area and is the most widely used semi-quantitative tool, but repeatability differs across infrared systems and inter-observer agreement is not guaranteed\cite{ngo2014repeat}. Two requirements follow: segmentation must support consistent area, length, and grading rather than visual appeal; and the eyelid margin must be segmented as well, otherwise the denominators of R\_loss and Meiboscore are unstable.

\subsection{Automated Segmentation of Meibomian Glands}

Early work used thresholds, morphology, and strip filters, and was sensitive to low contrast and device variation. Deep learning followed: Prabhu et al.\ segmented glands and quantified morphology\cite{prabhu2020segmentation}; Wang et al.\ quantified gland morphology and linked it to clinical grading\cite{wang2021morphology}; Liu et al.\ used scribble supervision to reduce pixel-level annotation cost\cite{liu2022scribble}; Setu et al.\ and Wang et al.\ explored deeper or lighter encoder--decoder variants\cite{setu2021deep,wang2019deep}; Yu et al., Zhang et al., and Lin et al.\ introduced density estimation, attention, and multi-scale features\cite{yu2022automatic,zhang2022density,lin2021automatic}; Liang et al.\ discussed cross-device and cross-dataset generalisation\cite{liang2024exploring}.

These studies establish feasibility, but their reporting differs substantially: whole-image versus sliding-window inference, Dice-only versus morphological metrics, and often unclear train/validation/test exclusivity. We do not dispute their conclusions; we place them under one protocol, with explicit tiling, threshold, and split signature.

\subsection{Public Benchmarks, Patient-Level Leakage, and Inference Confounders}

MGD-1K is the most representative public meibography dataset: 1000 infrared images, gland and eyelid masks, 320 subjects, and six expert Meiboscore rounds per image\cite{saha2022automated}. A 2025 multicentre study further shows that gland segmentation can reach high grading agreement at independent centres, but explicitly notes that its model covers only the upper eyelid and a single device\cite{li2025multicenter}; an MGD-1K baseline reports U-Net IoU around 0.72 and F1 around 0.82 without systematic patient-level splitting\cite{yu2025quantitative}. Public data are therefore available, but ``public'' does not imply ``comparable''.

Patient-level leakage is the easiest of these problems to hide at the implementation level. A subject in MGD-1K often has multiple images (upper/lower eyelid, left/right eye), and image-level random splits place different images of the same subject into training and test. Even with a patient-level split, directly using the official 200-case Meiboscore test list leaves many other images of those test patients in training. Few studies report both a file-list signature and an assignment signature, so readers cannot tell whether two experiments actually used the same split.

Tiling at inference is another hidden variable. Buglakova et al.\ show that stitching artefacts in sliding-window inference are tied to normalisation layers, and compare the trade-off between artifact-free and high-quality predictions\cite{buglakova2025tiling}. For 3--8-pixel glands, tile height also changes the effective receptive field and context ratio, and its effect can exceed architecture differences. We treat tiling as a first-class variable and build a multi-tile ensemble on top of it.

For uncertainty, the decision-aware protocol of Maganti and Gowda shows that TTA often gives a stronger error-detection signal than MC Dropout under a fixed review budget, but estimator rankings are not universal\cite{maganti2026uncertainty}; Geißler et al.\ show that structure size confounds the relationship between uncertainty and segmentation quality\cite{geissler2024structure}. In thin-structure tasks, the uncertainty source and structure size must therefore be reported together rather than as a single AUROC.

\subsection{Architectures for Medical Image Segmentation}

Encoder--decoder with skip connections remains a strong baseline\cite{ronneberger2015unet}; Attention U-Net gates irrelevant skip responses\cite{oktay2018attention}; UNet++ and MultiResUNet improve nested connectivity and multi-resolution processing\cite{zhou2018unetplusplus,ibtehaz2020multiresunet}; nnU-Net contributes a self-configured recipe---patch sampling, deep supervision, strong augmentation---rather than a module\cite{isensee2021nnunet}. Among Transformers, TransUNet uses a Transformer encoder\cite{transunet}, Swin-UNet builds a U-shaped window-attention network\cite{swinunet}, and SegFormer pairs a hierarchical MiT encoder with a lightweight head\cite{segformer}; UNETR and UNETR++ bring Transformers to 3D medical segmentation\cite{unetr,unetrpp}, and MedNeXt scales ConvNets with Transformer-inspired blocks\cite{mednext}. State-space models, represented by VM-UNet and U-Mamba, model long-range dependencies with linear complexity\cite{vmunet,umamba}.

For this task, long-range dependency is not the bottleneck; detail preservation is. SegFormer's finest stage has stride 4 and offers ADE20K/ImageNet pretraining, matching the 3--8-pixel target and the 999-case regime. Swin-UNet's windowed attention and four-stage downsampling are unfriendly to thin-line continuity; Mamba variants may fragment connectivity through 2D scan order and depend on compiled CUDA kernels, raising cross-platform reproduction cost. We therefore adopt SegFormer-B2 as the encoder and modify the decoder to retain 1/4 features; the other routes serve as literature comparisons and are not the default.

\subsection{Losses and Metrics for Thin Tubular Structures}

Region losses (Dice, cross-entropy) are sensitive to overlap but not to connectivity: when a gland breaks into two segments, Dice may drop only slightly while HD95 rises sharply. clDice constrains topology through soft-skeleton centerline overlap\cite{cldice}; skeleton-recall loss makes skeleton recall explicit\cite{skeletonrecall}; boundary loss and Boundary IoU focus optimisation and evaluation on the boundary band\cite{boundaryloss,boundaryiou}. Together they show that thin structures require multi-metric evaluation, since Dice alone hides topological failure.

We use clDice as the only skeleton-continuity term and do not add an explicit skeleton loss; the rationale is given in Section~\ref{sec:loss}: two terms pushing the same gradient direction amplify skeleton noise. Boundary loss remains an optional switch, disabled by default.

\subsection{Low-Contrast Enhancement, Domain Shift, and Inference Strategies}

CLAHE is the most common local contrast enhancement for infrared and medical images; its clip limit and window size control noise amplification and artefacts\cite{infraredclahe}. Deep low-light enhancement, such as Zero-DCE, offers a data-driven alternative\cite{zerodce,lowlightsurvey}, but usually targets visual quality rather than segmentation metrics. For domain shift, test-time adaptation and continual test-time adaptation adjust statistics without labels\cite{tent,cotta}; TTA reduces variance by averaging geometric and scale views\cite{wang2019aleatoric}.

We do not use test-time training, to avoid introducing adaptation on the test set; TTA is fixed to horizontal flips and two scales, and is enabled only for paper evaluation and final delivery. Cross-device variation is addressed mainly by intensity augmentation and a unified normalisation contract, not by modifying the model at test time.

\subsection{Interpretability, Uncertainty, and Reporting Standards}

Grad-CAM and its variants map gradients back to the input\cite{gradcam,gradcampp,scorecam}; uncertainty estimation separates aleatoric from epistemic, with TTA entropy and MC dropout as common approximations\cite{wang2019aleatoric,uncertaintyreview,phiseg}. Reporting standards include CLAIM, TRIPOD+AI, PROBAST, and REFORMS, covering imaging AI, prediction models, bias assessment, and reproducible research\cite{claim2024update,tripodai,probast,reforms}; metric selection itself has systematic recommendations and pitfalls\cite{metricsreloaded,metricpitfalls}.

We borrow their structure but operationalise the three items most relevant here: split signature, tiling-parameter reporting, and weight-status labelling. The system implements true Grad-CAM and TTA entropy maps, but we treat them as review aids rather than performance evidence.
